\section{Introduction}
\label{sec:intro}

Recent advances in multimodal foundation models~\citep{qwen3-vl,
Claude-Opus-4.6, Seed2.0} have strengthened visual understanding,
language reasoning, and tool-use capabilities~\citep{tu2026consensus}. Simultaneously, agent
systems~\citep{MM-MEM, milestone, Audio-Oscar} have rapidly evolved from
language only assistants toward interactive agents that can plan, invoke
tools, and operate in external environments. Graphical user interface
(GUI) agents~\citep{UI-TARS-2,Mobile-agent-v3.5} have therefore emerged
as a natural task form for connecting model intelligence with real digital
environments. GUI agents must understand screen content, plan operation
steps, and complete user goals through interface-level actions such as
clicking, typing, and swiping. Early studies largely centered on 
platform-specific web navigation or computer automation. Benchmarks such as OSWorld~\citep{osworld} and
MobileWorld~\citep{mobileworld} evaluate agents in interactive computer
and mobile environments, driving a shift from static interface
understanding toward long-horizon, platform-grounded interaction. Yet
real-world workflows often span computer applications, mobile apps, and web
services, requiring agents to adapt across heterogeneous GUI environments while preserving the interaction
conventions of each platform. This progression raises a central question:
\textit{how can a shared GUI agent continually adapt across heterogeneous
platforms while retaining platform-specific interaction behaviors}?

Recent efforts have attempted to extend GUI agents beyond isolated
platforms mainly by scaling interaction data or training a shared policy
on heterogeneous environments. Open-source datasets such as
OpenCUA~\citep{Opencua} and OpenMobile~\citep{openmobile} provide
growing collections of GUI trajectories, while multi-platform GUI
agents~\citep{Mobile-agent-v3.5, UI-venus-1.5} typically combine signals
from computer, mobile, or web environments through mixed supervised
fine-tuning (SFT), mixed reinforcement learning (RL), or model merging.
Although these efforts broaden platform coverage, they largely treat
cross-platform learning as aggregating more data or merging heterogeneous
training signals.

However, a capable cross-platform GUI agent requires more than exposure
to multiple platforms. It must acquire transferable interface reasoning
while preserving the distinct interaction conventions of each platform.
This creates two bottlenecks. First, high-quality cross-platform
trajectories remain scarce: existing datasets often focus on
single-platform settings and may contain invalid actions, inaccurate
state-action alignment, or inconsistent task granularity. Second,
computer and mobile platforms differ in action semantics and affordances;
for example, returning to the previous context may mean closing a window
on computers but pressing the back button in a mobile app. Naively
combining such signals can produce an overly averaged policy in joint
training and, under continual learning, cause catastrophic forgetting of
previously learned platform-specific behaviors.
Thus, cross-platform GUI learning requires stable platform-specific
behavioral anchors throughout policy optimization.

To address these challenges, we first build a unified cross-platform
data collection harness that collects interaction data from computer and
mobile environments with a consistent action interface and logging format.
Using this harness, we collect approximately 110K and 50K interaction
steps from computer and mobile environments, respectively. After rigorous
filtering and quality control, we construct \BENCH, a dataset
containing nearly 10K high-quality cross-platform interaction trajectories.
Building on Uni-GUI, we propose \METHODNAME, the first method to 
introduce multi-teacher on-policy distillation (MOPD) into GUI agent continual learning, 
with platform-conditioned teachers for multi-platform adaptation.


\METHODNAME dynamically selects a platform-specific teacher for each
rollout environment and transfers platform-specific behavior distributions
to a shared policy. For computer and mobile environments, the model aligns
with native interaction patterns preserved by corresponding teachers,
preventing behavior signals with different interaction conventions from
being indiscriminately mixed during optimization. These platform-specific
teachers provide more than additional supervision. They serve as stable
behavioral anchors, allowing a shared policy to improve task completion
while preserving platform-specific behavioral priors. Through such
environment-conditioned policy alignment, \METHODNAME better balances
adaptation to new platforms with retention of existing platform behaviors,
mitigates behavioral convention mixing and catastrophic forgetting of
platform-specific behaviors, and encourages a
platform-aware policy that activates different interaction modes according
to environment context.



Empirical results further demonstrate the effectiveness of \METHODNAME. It achieves
task success rates of $\mathbf{38.2\%}$ on OSWorld and $\mathbf{12.0\%}$ on MobileWorld, corresponding to
relative improvements of $\mathbf{12.7\%}$ and $\mathbf{55.8\%}$ over the base model, respectively.
Importantly, these gains are observed in both computer and mobile environments,
suggesting that \METHODNAME improves adaptation to new platforms without sacrificing
existing platform capabilities and enables more balanced continual
optimization across platforms.


The main contributions of this work are summarized as follows:

\begin{itemize}
    \item We introduce \BENCHhl, a high-quality dataset for multi-platform GUI interaction. Using a unified cross-platform data-collection harness, we collect approximately 10K high-quality cross-platform interaction trajectories from computer and mobile environments.

    \item We propose \METHODNAMEhl, the first method to introduce multi-teacher on-policy distillation for continual GUI agent learning, addressing behavioral convention mixing, platform-specific degradation, and catastrophic forgetting in continual GUI agent learning.

    \item We conduct systematic experiments on \textbf{OSWorld} and \textbf{MobileWorld}. \METHODNAMEhl achieves task success rates of $\mathbf{38.2\%}$ and $\mathbf{12.0\%}$ on OSWorld and MobileWorld, respectively, demonstrating its effectiveness in preserving cross-platform capabilities while adapting to new platforms.
\end{itemize}
